\subsection{形式幂级数集合上的拓扑。可和族}

根据定义，A{[}{[}I{]}{]} 无非就是乘积集合 \(\mathbf{A}^{\mathbf{N}^{(I)}}\) 。除非另有明确说明，我们将赋予 A 离散拓扑，并赋予 A{[}{[}I{]}{]} 乘积拓扑（一般拓扑学，I，第 31 页及以下），我们称之为典范拓扑。赋有加法与离散拓扑的 A 是一个分离且完备的拓扑群；因此就加法而言，A{[}{[}I{]}{]} 是一个分离且完备的拓扑群（一般拓扑学，III，第 238 与 242 页，以及一般拓扑学，II，第 187 页）。此外，多项式代数 A{[}(X \(_{i}\) ) \(_{i\in I}\) {]} 在 A{[}{[}I{]}{]} 中稠密（一般拓扑学，III，第 238 页，命题 25），因此我们可以把 A{[}{[}I{]}{]} 视为 A{[}(X \(_{i}\) ) \(_{i\in I}\) {]} 的完备化。

对每个 \(\beta \in \mathbf{N}^{(I)}\) ，令 \(S_{\beta}\) 为多重指标 \(\nu\) （满足 \(\nu \leqslant \beta\) ）组成的集合，并令 \(\alpha_{\beta}\) 为由所有形式幂级数 \(u = \sum_{\nu} \alpha_{\nu} X^{\nu}\) （满足 \(\alpha_{\nu} = 0\) ，对

\(\nu \in S_{\beta}\) ）组成的集合。显然 \(S_{\beta}\) 是 \(\mathbf{N}^{(I)}\) 的有限子集，而 \(\mathbf{N}^{(I)}\) 的每个有限子集都包含于形如 \(S_{\beta}\) 的集合。由此可见，族 \((\mathfrak{a}_{\beta})_{\beta \in \mathbf{N}^{(I)}}\) 是 A{[}{[}I{]}{]} 中 0 的一个基本邻域系。集合 \(\mathfrak{a}_{\beta}\) 是 A{[}{[}I{]}{]} 中的理想，因此（一般拓扑学，III，第 275 页）A{[}{[}I{]}{]} 是拓扑环。

引理 1. --- 设 L 是一个无限集， \((u_{\lambda})_{\lambda \in L}\) 是 A{[}{[}I{]}{]} 的元素的一个族，并令 \(u_{\lambda} = \sum_{\nu} \alpha_{\lambda, \nu} X^{\nu}\) （对 \(\lambda \in L\) ）。则以下条件是

等价的：

\begin{enumerate}
\def\labelenumi{(\roman{enumi})}
\item
  族 \((u_{\lambda})_{\lambda \in L}\) 在 A{[}{[}I{]}{]} 中是可和的（一般拓扑学，III，第 262 页）。
\item
  我们有 \(\lim u_{\lambda} = 0\) ，这是沿 \(L\) 的有限子集的补集组成的滤子取的极限。
\item
  对每个 \(\nu \in \mathbf{N}^{(I)}\) ，我们有 \(\alpha_{\lambda, \nu} = 0\) ，除去有限个指标 \(\lambda \in L\) 外。
\end{enumerate}

当这些条件成立时，级数 \(u = \sum_{\lambda \in L} u_{\lambda}\) 等于 \(\sum_{\nu} \alpha_{\nu} X^{\nu}\) ，其中 \(\alpha_{\nu} = \sum_{\lambda \in L} \alpha_{\lambda, \nu}\) （对每个 \(\nu \in \mathbf{N}^{(I)}\) ）。

（i）与（ii）的等价性由一般拓扑学 III，第 263 页，推论 2 得出。（ii）与（iii）的等价性由乘积空间中极限的性质得出（一般拓扑学，I，第 55 页，推论 1）。

最后一个断言由《一般拓扑学》III，第266页，命题4推出。

下面给出一些可和族的例子。

\begin{enumerate}
\def\labelenumi{\alph{enumi})}
\item
  设 \(u \in \mathbf{A}[[\mathrm{I}]]\) ，并令 \(\alpha_{\nu}\) 是 \(\mathbf{X}^{\nu}\) 在 \(u\) 中的系数。则族 \((\alpha_{\nu}\mathbf{X}^{\nu})_{\nu \in \mathbb{N}^{(I)}}\) 是可和的，其和为 \(u\) （这证明了写 \(u = \sum_{\nu} \alpha_{\nu}\mathbf{X}^{\nu}\) 的合理性）。
\item
  设 \(u \in \mathbf{A}[[\mathbf{I}]]\) ；对每个整数 \(p \geqslant 0\) ，令 \(u_p\) 为次数 \(p\) 的齐次分量（ \(u\) 的）。则族 \((u_p)_{p \in \mathbb{N}}\) 是可和的，并且我们有 \(u = \sum_{p \geqslant 0} u_p\) 。
\item
  设 \((u_{\lambda})_{\lambda \in L}\) 是 A{[}{[}I{]}{]} 的元素的一个族，并假设对每个整数 \(n \geqslant 0\) ，由所有 \(\lambda \in L\) （满足 \(\omega(u_{\lambda}) < n\) 者）组成的集合是有限的。则族 \((u_{\lambda})_{\lambda \in L}\) 是可和的。
\end{enumerate}

注记. --- 假设 I 是有限的。对每个整数 \(n \geqslant 0\) ，令 \(\mathbf{b}_n\) 为由所有形式幂级数 \(u \in \mathbf{A}[[\mathbf{I}]]\) （满足 \(\omega(u) \geqslant n\) 者）组成的集合。序列 \((\mathbf{b}_n)_{n \geqslant 0}\) 是 A{[}{[}I{]}{]} 中 0 的一个基本邻域系。因此，由元素 \(u_{\lambda}\) （属于 A{[}{[}I{]}{]}， \(\lambda \in L\) ）组成的族是可和的，当且仅当对每个 \(n \in \mathbb{N}\) ，由所有 \(\lambda \in L\) （满足 \(\omega(u_{\lambda}) < n\) 者）组成的集合是有限的。

命题 1. --- 设 \((u_{\lambda})_{\lambda \in L}\) 与 \((v_{\mu})_{\mu \in M}\) 是 A{[}{[}I{]}{]} 的元素的两个可和族。则族 \((u_{\lambda}v_{\mu})_{(\lambda, \mu) \in L \times M}\) 是可和的，并且我们有

\[
\sum_ {(\lambda , \mu) \in L \times M} u _ {\lambda} v _ {\mu} = \left(\sum_ {\lambda \in L} u _ {\lambda}\right) \left(\sum_ {\mu \in M} v _ {\mu}\right).\tag{2}
\]

设 \((\alpha_{\lambda, \nu})_{\nu \in \mathbf{N}^{(I)}}\) （相应地 \((\beta_{\mu, \nu})_{\nu \in \mathbf{N}^{(I)}}\) ）是 \(u_{\lambda}\) （相应地 \(v_{\mu}\) ）的系数族。对每个 \(\nu \in \mathbf{N}^{(I)}\) ，只存在有限多对 \((\nu_1, \nu_2) \in \mathbf{N}^{(I)} \times \mathbf{N}^{(I)}\) 满足 \(\nu_1 + \nu_2 = \nu\) ，因此只有有限多对 \((\lambda, \mu) \in L \times M\) 使得 \(X^{\nu}\) 在 \(u_{\lambda}v_{\mu}\) 中的系数 \(\neq 0\) 。因此族 \((u_{\lambda}v_{\mu})_{(\lambda, \mu) \in L \times M}\) 是可和的。而公式 (2) 由和的结合性得出（一般拓扑学，III，第 265 页，公式 (2)）。

在 A{[}{[}I{]}{]} 中，乘积是一个结合且交换的合成法则。因此我们可以谈论 A{[}{[}I{]}{]} 的元素的可乘族以及一个可乘族的乘积（一般拓扑学，III，第 262 页，注记 3）。

命题 2. --- 设 \((u_{\lambda})_{\lambda \in L}\) 是 A{[}{[}I{]}{]} 的元素的一个可和族。(i) 族 \((1 + u_{\lambda})_{\lambda \in L}\) 是可乘的。

\begin{enumerate}
\def\labelenumi{(\roman{enumi})}
\setcounter{enumi}{1}
\tightlist
\item
  设 \(\mathfrak{F}\) 是 L 的所有有限子集组成的集合。对任意 \(\mathbf{M} \in \mathfrak{F}\) ，令 \(u_{\mathbf{M}} = \prod_{\lambda \in \mathbf{M}} u_{\lambda}\) 。则族 \((u_{\mathbf{M}})_{\mathbf{M} \in \mathfrak{F}}\) 是可和的，并且我们有
\end{enumerate}

\[
\sum_ {M \in \mathfrak {F}} u _ {M} = \prod_ {\lambda \in L} (1 + u _ {\lambda}).
\]

按本节开头那样定义理想 \(\mathfrak{a}_{\beta}\) ，并设 \(\beta \in \mathbf{N}^{(I)}\) 。存在一个有限子集 \(L_0\) （ \(L\) 的），使得 \(u_{\lambda} \in \mathfrak{a}_{\beta}\) （对 \(\lambda \notin L_0\) ）。于是对每个 \(M \in \mathfrak{F}\) （满足 \(M \subset L_0\) ），我们有 \(u_M \in \mathfrak{a}_{\beta}\) 。由此可见，族 \((u_M)_{M \in \mathfrak{F}}\) 是可和的。另一方面，对任意有限子集 \(M_0\) （ \(L\) 的），我们有

\[
\sum_ {M \subset M _ {0}} u _ {M} = \prod_ {\lambda \in M _ {0}} (1 + u _ {\lambda}).
\]

沿滤过有序集 \(\mathfrak{F}\) 取极限，左端的极限为 \(\sum_{M\in\mathfrak{F}}u_{M}\) 。因此右端的极限也为 \(\sum_{M\in\mathfrak{F}}u_{M}\) ，这就证明了 (i) 与 (ii)

同时。

命题 3. --- 设 \(u = \sum_{\nu} \alpha_{\nu} X^{\nu} \in \mathbf{A}[[\mathbf{I}]]\) ，并设 \(m\) 是一个 \(> 0\) 的整数。对每个 \(n \in \mathbb{N}\) ，令 \((\alpha_{\nu, n})_{\nu \in \mathbb{N}^{(I)}}\) 是 \(u^n\) 的系数族。若 \(\alpha_0^m = 0\) ，则 \(\alpha_{\nu, n} = 0\) （对 \(n \geqslant |\nu| + m\) ）。

设 \(\nu \in \mathbf{N}^{(I)}\) 与 \(n\in \mathbf{N}\) 。我们有

\[
\alpha_ {\nu , n} = \sum_ {\nu (1) + \dots + \nu (n) = \nu} \alpha_ {\nu (1)} \dots \alpha_ {\nu (n)}.
\]

如果 \(n \geqslant |\nu| + m\) 且 \(\nu(1) + \cdots + \nu(n) = \nu\) ，我们有 \(|\nu(1)| + \cdots + |\nu(n)| \leqslant n - m\) 。于是 \(\nu(r) = 0\) 对至少 m 个不同的 r 值成立，从而 \(\alpha_{\nu(r)} = \alpha_0\) ；由此推出 \(\alpha_{\nu(1)} \ldots \alpha_{\nu(n)} = 0\) ，从而得到结论。

推论. --- 设 \(u \in \mathbf{A}[[\mathbf{I}]]\) ；那么要使 \(\lim_{n \to \infty} u^n = 0\) 成立，其必要且

充分条件是 \(u\) 的常数项为幂零元。

设 \(\alpha_0\) 是 \(u\) 的常数项。 \(u^n\) 的常数项是 \(\alpha_0^n\) ，因此所述条件是必要的；其充分性由命题 3 得出。
% semantic-fragments: legacy-input-seam
\relax{}
